\section{20 декабря 1917 г.}

\lp{20260603_193855019}

Цюрих
Письмо 115-ое

С *** констатировать
сейчас, что мое последнее письмо 114-ое отправлено 28 ноября, то есть больше
трех недель тому назад. После я отправил Вам лишь открытку с *** скоро написать.
Я не знаю, что делается со мною, но я *** не *** если бы мне сказали, что
я Вам не писал лишь немного дольше ***
Простите моя самая дорогая!
И не подумайте, что я не хотел писать, или, что я ***

Сегодня день моего рождения. Очень тяжело. Тяжело очень. ***
Одна надежда на бога.

\lp{20260603_193915733}

Вспомнил, что два года тому
назад Вы прислали мне
в письме Ваше любимое
стихотворение Лермонтова,
Ваш лейтмотив\textellipsis{} Из
него: *** ты о них (о *** и об облаках лишь вспомяни, будь к
земному без ***
как они.

Так хочется сказать своей душе\textellipsis{}
*** плохо спал: все думал о жизни
и о\textellipsis{} смерти; не в смысле
страха смерти, но в смысле:
смерти страха; ибо если
страх перед богом есть
начало всякой премудрости,
то страх перед ***
есть начало
всякой немудрости; и страх
перед жизнью еще больше
опустошителен, нежели

\lp{20260603_193938558}

\noindent
страх перед смертью.
И гораздо большее число
заражено страхом жизни
нежели страхом смерти!
Только последнее считается
особенно для мужчины позором
и потому об этом
говорят; тогда как страх
жизни недостаточно осознан
как таковой, ***
в нашей природе,
проявляется крайне
различно, и поэтому о нем
к сожалению не говорят.
Страх жизни ***,
нежели страx смерти.
Самоубийцы это превосходно
*** и наглядно.
Георгиевские кавалеры
способствуют доказательству
того же положения.
Многие не убивают
себя лишь потому, что
одержимы обоими страхами


\lp{20260603_193950308}

\noindent
в одинаковой мере.
О них можно сказать,
что они живут
(то есть не умирают) потому,
что *** -- Но этот
плюс, получаемый от двух
минусов, есть небытие,
притворившееся бытием.

Разсуждения для дня рождения
недостаточно торжественно
веселые, но зато очень правдивые.
Завтра напишу больше
и радостнее. Хочу это
отправить сегодня
[redacted] *** моего рождения
сегодня утром, когда я пил кофе,
пришел Goethe-Kalender 1918
из книжного магазина. ***
рад этому подарку, преподнесенному
мне ***аем.
Прилож*** портрет Geor***,
*** смотрит в глаза
зрителю.

Обнимаю Колю.

Ваш Миля.
